\bookchapter{第五章：使手機更知汝意——基礎設置}

\chapterintro{花數分設，體驗大升。此等設置猶爲手機“調音”，使更合汝用之習。}

\booksection{一、聲與震}

設一喜之鈴聲與通知聲。活用靜音之模（開會）、勿擾之模（寢時），世界頓清。

\begin{itemize}
  \item \textbf{鈴聲}：擇一汝能聞、不尬者。
  \item \textbf{震之模}：靜音時賴震以醒。
  \item \textbf{勿擾之模}：可設例外，如允通訊錄來電。
  \item \textbf{音量之離}：媒體之量、鈴聲之量、鬧鐘之量可分調。
\end{itemize}

\booksection{二、顯與亮}

\begin{itemize}
  \item \textbf{自動之亮}：使手機據環境之光自調，省心。
  \item \textbf{護眼/夜覽之模}：濾藍光，夜間觀手機舒。
  \item \textbf{深色之模}：省電（于OLED屏尤效），且夜間不刺眼。諸多App亦支持深色之模。
  \item \textbf{屏刷新之率}：部分安卓機可調，越高（如120Hz）滑越絲滑。
  \item \textbf{自動鎖屏}：設十五秒或三十秒，安且省電。
  \item \textbf{字體之大}：覺字小？調大可也。
\end{itemize}

\booksection{三、壁紙、主題與桌}

換一佳之壁紙（靜態或動態），心情皆變。安卓品牌手機常有主題之店，能整套換圖標與風。

\textbf{桌面小組件}：天氣、日程、音樂之控置桌面，一步直達。iOS與安卓皆支持，長按桌面空處加之。

\booksection{四、效率之設：使手機爲汝之“外掛腦”}

此部分乃“一設則難返”之功，每皆能使汝日常之效大升。

\booksubsection{專注/免打擾之模：奪被通知所奪之注}

\begin{itemize}
  \item \textbf{核心之思}：非"少觀手機"，乃"當專注時專注，當聯時再聯"。
  \item \textbf{用之建議}：\newline · 工作之模：惟允"星標"聯繫人與必要App過；\newline · 睡眠之模：自動開，惟允鬧鐘；\newline · 健身/閱/會議：據當前活動定製。
  \item \textbf{進階之玩}：用自動化加位置/時觸，使模自切。如：iOS 快捷指令可設"至公司即自動入工作之模"。
\end{itemize}

\booksubsection{通知之管：使紅點不“奪命連環 call”}

\begin{itemize}
  \item \textbf{閉營銷通知}：九成之推實汝不需。
  \item \textbf{按要之序}：將人際通訊之類（電、短信、即時通訊）置頂；新聞、購物、社交平臺之類可設"靜默推"。
  \item \textbf{分組通知}：將同類通知折於一，避一解鎖被數十條淹沒。
\end{itemize}

\booksubsection{快捷之手：使操快人一步}

\begin{itemize}
  \item \textbf{截屏}：電源加音量下（iOS全面屏）/電源加音量下（安卓通用）。
  \item \textbf{截長圖/滾動截屏}：常截後點"長截圖"之擇。
  \item \textbf{一鍵掃碼}：控制中心/快捷設置中加"掃一掃"之圖標。
  \item \textbf{輕點背面（iPhone）}：於 [設置] > [輔助功能] > [觸控] > [輕點背面] 中綁"截屏"或"鎖屏"等動作。
  \item \textbf{語音之喚}：Siri/小愛同學/小藝/小布，喊一聲即能查天氣、設鬧鐘。
\end{itemize}

\booksubsection{輸入法：最被低估之效神器}

\begin{itemize}
  \item \textbf{擇對輸入法}：\newline · 搜狗/百度：詞庫大，流行語更速；\newline · 訊飛：語音輸入識別之率至高，開會記筆記之神器；\newline · Gboard（谷歌）：簡、跨平臺同步強、支持滑行輸入；\newline · iOS 自帶：隱私至佳，與系統深集成。
  \item \textbf{必備之開關}：智能糾錯、聯想詞、雲詞庫同步、上滑輸入數字符號。
  \item \textbf{隱私之提}：慎用"整句聯想"或"個性化推"之功，或上輸輸入內容。
\end{itemize}

\booksubsection{桌與多任：使工作流順暢}

\begin{itemize}
  \item \textbf{App 分屏/畫中畫}：觀視頻時回消息，上網時記筆記。
  \item \textbf{小窗/懸浮窗}：微信/抖音支持"小窗之模"，回覆畢即收。
  \item \textbf{常用 App 置頂}：Dock欄（屏底）惟放最高頻之四，減查尋之時。
\end{itemize}

\begin{tcbinfo}{"上工第一日"檢之單}
“上工第一日”檢之單\newline (1) 鎖屏密碼加生物識 (2) 通知分組加閉營銷推 (3) 字體之大/刷新之率/深色之模\newline (4) 默認輸入法/瀏覽器/郵件 (5) 緊聯人/醫療卡 (6) 備份設置

\end{tcbinfo}

\booksection{五、輔助功能與無障：科技向善之力}

莫爲“輔助功能”之名所嚇，此等功能人人能用。其初爲殘障者設，然於常人亦有用——臨手傷抱娃、地鐵上耳機壞不聞、戴老花鏡觀小字……皆能救場。說到底，無障之設體現了\textbf{科技向善}之理念：使每個人皆能平等地享數字之生活。

\booksubsection{視障用戶之“眼”——屏閱讀器全攻略}

屏閱讀器乃視障用戶用手機之核心之具。其將屏上\textbf{所有可視內容朗而出}，用戶經手與語音以操手機。下分別介 iOS “旁白”與安卓 TalkBack 之用法。

\begin{tcbinfo}{iOS 旁白（VoiceOver）}
\textbf{開啓之徑}：[設置] > [輔助功能] > [旁白]，開開關。亦可三擊側邊鍵（或主屏鍵）速開。

\textbf{基本之手}：

\begin{enumerate}
  \item 單指滑：逐元素覽屏內容；
  \item 單指滑：逐元素覽屏內容；
  \item 雙指輕掃：上下/左右整頁滾；
  \item 三指輕掃：切頁；
  \item 三指輕掃：切頁；
\end{enumerate}

\textbf{個性化之設}：[旁白] > [語音]，可調\textbf{朗之速}（默認約五成，可調至七至八成聽感更流）、\textbf{音調}、切不同語音包。iOS自帶“增強語音”與“語音控制器”，前者能減呼吸聲與停頓，聽感更自然。

\end{tcbinfo}

\begin{tcbinfo}{、切不同語音包。iOS自帶"增強語音"與"語音控制器"，前者能減呼吸聲與停頓，聽感更自然。}
\textbf{開啓之徑}：[設置] > [輔助功能] > [TalkBack]，開開關。或長按音量上下鍵三秒速啓（部分品牌支持）。

\textbf{基本之手}：

\begin{enumerate}
  \item 單指左右滑：切焦點之元素；
  \item 單指左右滑：切焦點之元素；
  \item 單指雙擊：激活元素；
  \item 雙指上下/左右滑：滾；
  \item 四指單擊：暫停/續朗。
\end{enumerate}

\textbf{個性化之設}：[TalkBack] > [文字轉語音設置]，可調\textbf{語速}、\textbf{音調}、切\textbf{語音引擎}（建議裝 Google 文字轉語音引擎，朗更自然）。

\end{tcbinfo}

\booksubsection{使小字變大、使模糊變清——視覺之調}

\booksubsection{放大之手}

\begin{itemize}
  \item \textbf{iOS}：[設置] > [輔助功能] > [縮放]，開後\textbf{三指輕點屏}即放大。默認全屏縮，可擇窗口縮（局部放大窗）。\textbf{雙指開合可調倍（1.5x 至 15x）}，三指雙擊鎖當前縮之比。
  \item \textbf{安卓}：[設置] > [輔助功能] > [放大]，三指點擊放大，再三指點擊退；或直\textbf{三指點擊屏三下}啓臨放大（最多八倍）。
\end{itemize}

\booksubsection{高對比之字/粗體文本}

\begin{itemize}
  \item \textbf{iOS}：[設置] > [輔助功能] > [顯示與文字大小]，開"粗體文本"、"增強對比度"、"降低透明度"等開關。"智能反轉"可於淺色底使文字保持黑，避白底白字。
  \item \textbf{安卓}：[設置] > [顯示] > [字體大小]，可調"小、普通、大、超大"四檔。華爲/小米額外供"粗體"與"高對比文本"之開關。
\end{itemize}

\booksubsection{色彩濾鏡（色盲之友）}

\begin{itemize}
  \item \textbf{iOS}：[設置] > [輔助功能] > [顯示與文字大小] > [色彩濾鏡]，供\textbf{紅/綠/藍/黃四色盲之模}，及"灰度"與"反轉顏色"之擇。建議先做一次色覺之測（系統自帶）再擇對應濾鏡。
  \item \textbf{安卓}：[設置] > [輔助功能] > [顏色校正]，含\textbf{三色盲之模}加"反轉顏色"之開關。
\end{itemize}

\booksubsection{AI 圖像之述：使AI爲汝之“眼”}

當屏閱讀器遇圖時，傳統之做法乃跳或惟讀“圖”二字。然今AI可\textbf{以一語述圖之內容}，於識驗證碼、讀圖、識菜單甚有助。

\begin{tcbinfo}{諸系統 AI 圖像之述支持之況}
\begin{itemize}
  \item \textbf{iOS 15+}：於 [設置] > [輔助功能] > [旁白] > [語音描述] 中開"圖像之述"。系統將調\textbf{本地 AI 之模}自動生圖之文述（須 iPhone 12 及以上）。
  \item \textbf{新版 TalkBack 應用}：Google TalkBack 于較新版中集\textbf{圖之問}之功，可問"此圖中有何"，AI 實時答（不嚴賴安卓系統之版，更 TalkBack 應用即能體驗）。
  \item \textbf{第三方 App}："Seeing AI"（iOS）、"Lookout"（安卓）、"小藝幫幫"（華爲）皆爲優之AI視覺之助，可識物、讀文件、識人臉之表情。
\end{itemize}

\end{tcbinfo}

\begin{tcbpractice}{使用之場}
使用之場\textbf{此乃真之“科技使生活更平等”}。

\end{tcbpractice}

\booksubsection{聽障用戶之“耳”——實時字幕全攻略}

實時字幕（Live Captions）乃\textbf{系統級之語音轉文字}之功，聽障用戶於噪之環、圖書館、外語之場皆特有用。

\begin{tcbinfo}{實時字幕之開與優化}
\textbf{iOS}：[設置] > [輔助功能] > [實時字幕]，開開關。系統將檢設備所有音頻（通話、視頻、播客、FaceTime、語音備忘），自動顯字幕於屏上。

\textbf{安卓 12+}：按\textbf{音量鍵}點擊彈出框中之“實時字幕”圖標；或於 [設置] > [輔助功能] > [實時字幕] 中開。

\textbf{準確率優化之巧}：

\begin{enumerate}
  \item 盡量于\textbf{靜之環}用，遠離空調、機之轟鳴；
  \item 言者\textbf{咬字清、語速適}時識別之率至高；
  \item 時識別之率至高；
  \item 支持中文普話、粵語、英語等主流之種，\textbf{小語種識別之率較低}；
  \item ；
\end{enumerate}

\end{tcbinfo}

\booksubsection{通話之場雙道：語音轉字加字轉音}

當用戶\textbf{既不聞亦不能言}時（如突失聲、喉術恢復期、重口喫），手機猶能當“傳話員”。

\begin{tcbinfo}{通話之實時轉譯（TTY / RTT）}
\textbf{iOS}：[設置] > [輔助功能] > [RTT/TTY]，開後通話時屏上顯鍵盤，對方之音被實時打成字，\textbf{汝所打之字亦被念出}予對方聽。

\textbf{安卓}：[設置] > [輔助功能] > [實時轉錄] 加 [聲音通知]。“實時轉錄”將對方言實時顯爲字；“聲音通知”於不便語音時將汝之入轉爲語音放。

\textbf{技術路線}：

\begin{itemize}
  \item \textbf{RTT（Real-Time Text）}：雙方皆須支持，字邊打邊顯，無延遲，主流運營商已支持；
  \item \textbf{TTY（電傳打字）}：須專用TTY設備，已漸爲RTT所代；
  \item 微信/QQ之\textbf{語音轉字之功}亦爲替之方：長按語音消息 "轉字"。
\end{itemize}

\end{tcbinfo}

\booksubsection{向所有人之“小而實用”之功}

\begin{itemize}
  \item \textbf{輕點背面}（iOS 14+）：[設置] > [輔助功能] > [觸控] > [輕點背面]，可綁\textbf{雙擊/三擊}觸截圖、鎖屏、控制中心等動作。雙手持快遞時特好用。
  \item \textbf{輔助之觸/懸浮球}：iOS [輔助功能] > [觸控] > [輔助之觸]，安卓 [輔助功能] > [懸浮導航]。虛按鈕替手或物理之鍵，鍵不靈或單手操時救場。
  \item \textbf{切換之控}：經外部開關、單次點擊、頭部動以控手機，\textbf{爲肢體障者設}。iOS [輔助功能] > [切換之控]，安卓 [輔助功能] > [開關之控]。
  \item \textbf{聲之識}（iOS）：[輔助功能] > [聲之識]，開後手機將"聞"門鈴、嬰哭、火警、警笛等聲，\textbf{識到則震加屏通知}，於聽障用戶與睡期間甚有值。
  \item \textbf{LED 閃之醒}：iOS [輔助功能] > [音頻/視覺] > [LED 閃以醒]，安卓 [設置] > [通知] > [閃光燈提醒]。來電與通知時閃光燈閃，\textbf{噪之環無聲亦能感}。
  \item \textbf{單聲道之音}：[設置] > [輔助功能] > [音頻] > [單聲道之音]，將左右聲道合。單耳聽力損者不失另一半聲之信。
\end{itemize}

\begin{tcbinfo}{蘋果之無障設之哲學}
蘋果之無障設之哲學

\end{tcbinfo}

\begin{tcbpractice}{花三十秒略體驗之}
花三十秒略體驗之

\begin{enumerate}
  \item 開\textbf{旁白/TalkBack}，試不賴眼用三指滑以覽微信；
  \item 開\textbf{實時字幕}，對手機言以觀識別之準率；
  \item ，對手機言以觀識別之準率；\textbf{三指點擊}放大一段字至二百\%，體驗視障者之日用；
  \item 放大一段字至二百\%，體驗視障者之日用；
\end{enumerate}

\end{tcbpractice}
